\begin{otherlanguage}{english}
\noindent This project measures the extent to which the third-party tracking software development kits (SDKs) embedded in Brazilian government Android applications are disclosed in their privacy policies and in Google Play's Data safety section, using as the unit of analysis the pair formed by an application and each SDK found in it. Based on studies that found missing or incorrect disclosures in a relevant share of the applications they analysed, the hypothesis is that at least 10\% of app--SDK pairs will lack a consistent disclosure. The design is quantitative and observational. The current version of each application will be downloaded from Google Play and identified by its file hash; SDKs will be identified by three detectors, whose agreement will be measured with Fleiss' $\kappa$; each app--SDK pair will be classified according to what the documents disclose; and the network traffic of a subsample will be observed in an emulator. An SDK's presence in application code will not be treated as proof of data transmission or of a violation of the Brazilian General Data Protection Law (LGPD). The study is limited to Android, which runs most smartphones in Brazil and whose packages can be obtained and analysed with open tools. The preparatory search identified 396 candidate government applications, 376 of which were preliminarily classified as official. Because AndroZoo held only 92 of the 396 candidates, every application will be downloaded from Google Play. In a 20-package pilot, the share of applications with at least one tracker ranged from 75\% to 95\%, depending on the detector. The systematic review identified 396 records published from 2018 to 2026 with topic terms in the title, screened 312 after removing duplicates, and forwarded 148 for full-text assessment, with 31 inclusion proposals so far. Main data collection and hypothesis testing are planned for Capstone Project II.

\end{otherlanguage}
